\subsection{Triangularization and Decomposition}

\begin{theorem} \label{thm:complex-operator-has-eigenvalue} \hlt{Complex Operators Have Eigenvalues}\\
Every operator on a finite-dimensional, nonzero complex vector space has an eigenvalue.
\begin{proof}
Let $T\in\LL(\mathcal{V})$. By Proposition \ref{prop:finite-dimensional-operators-have-annihilating-polynomials}, $T$ has a nonconstant annihilating polynomial. Since $\mathcal{V}$ is complex and nonzero, Theorem \ref{thm:annihilating-polynomials-control-spectrum} implies that one root of this polynomial is an eigenvalue of $T$.\\
\hlt{Structural Commentary}: Finite dimension gives an annihilating polynomial; over $\mathbb{C}$, it factors linearly.\\
\hlt{Key Insight}: One noninjective factor produces an eigenvector.
\end{proof}
\end{theorem}

\begin{definition} \label{def:matrix-commutant} \hlt{Commutant}\\
For $\emptyset\neq\mathcal{F}\subseteq\operatorname{Mat}(n,n,\mathbb{F})$, the \hlt{commutant} of $\mathcal{F}$ is $\operatorname{Comm}(\mathcal{F})\defeq\{X\in\operatorname{Mat}(n,n,\mathbb{F}): XA=AX\ \forall A\in\mathcal{F}\}$.\\
For a single matrix $A$, write $\operatorname{Comm}(A)$.
\end{definition}

\begin{lemma} \label{lem:block-diagonal-commutant-distinct-scalars} \hlt{Commuting with a Block Scalar Diagonal Matrix}\\
Let $\lambda_1,\ldots,\lambda_k\in\mathbb{F}$ and let $\Lambda=\lambda_1 I_{n_1}\oplus\cdots\oplus\lambda_k I_{n_k}$.\\
Partition $A=[A_{ij}]\in\operatorname{Mat}(n,n,\mathbb{F})$ conformally with $\Lambda$, so $A_{ij}\in\operatorname{Mat}(n_i,n_j,\mathbb{F})$.
\begin{enumerate}[label=(\roman*)]
\setlength{\itemsep}{0pt}
\item If $A$ is block diagonal, then $A\Lambda=\Lambda A$.
\item If $A\Lambda=\Lambda A$ and $\lambda_1,\ldots,\lambda_k$ are distinct, then $A$ is block diagonal.
\end{enumerate}
\begin{proof}
Block multiplication gives $(A\Lambda)_{ij}=\lambda_jA_{ij}$ and $(\Lambda A)_{ij}=\lambda_iA_{ij}$. If $A$ is block diagonal, then $A_{ij}=0$ for $i\neq j$, and these two expressions agree for every block. Conversely, if $A\Lambda=\Lambda A$, then $(\lambda_j-\lambda_i)A_{ij}=0$ for every $i,j$. Distinctness of the $\lambda_i$ forces $A_{ij}=0$ whenever $i\neq j$, so $A$ is block diagonal.\\
\hlt{Structural Commentary}: Distinct scalar blocks force commuting matrices to respect the block decomposition.\\
\hlt{Key Insight}: Distinct block scalars make every off-diagonal block visible and force it to vanish.
\end{proof}
\end{lemma}

\begin{example} \hlt{Distinctness Is Necessary}\\
The distinctness hypothesis in Lemma \ref{lem:block-diagonal-commutant-distinct-scalars} cannot be dropped. If $\lambda_1=\cdots=\lambda_k=1$, then $\Lambda=I_n$ commutes with every matrix in $\operatorname{Mat}(n,n,\mathbb{F})$, including matrices that are not block diagonal for the chosen partition.
\end{example}

\begin{theorem} \label{thm:commuting-family-common-eigenvector} \hlt{Common Eigenvector for a Commuting Family}\\
Let $\mathcal{V}\neq\{\mathbf{0}\}$ be finite-dimensional over $\mathbb{C}$, and let $\mathcal{F}\subseteq\LL(\mathcal{V})$ be a nonempty commuting family: $ST=TS$ for all $S,T\in\mathcal{F}$. Then there exists $\mathbf{v}\neq\mathbf{0}$ that is an eigenvector of every operator in $\mathcal{F}$.
\begin{proof} First suppose $\mathcal{F}=\{T_1,\ldots,T_m\}$ is finite, and induct on $m$. The case $m=1$ is Theorem \ref{thm:complex-operator-has-eigenvalue}.\\
For $m>1$, choose $\lambda\in\operatorname{spec}(T_1)$ and set $\mathcal{U}=E_{\lambda}(T_1)$. Then $\mathcal{U}\neq\{\mathbf{0}\}$.\\
If $\mathbf{u}\in\mathcal{U}$ and $j\geq2$, commutation gives $(T_1-\lambda I_{\mathcal{V}})T_j\mathbf{u} =T_j(T_1-\lambda I_{\mathcal{V}})\mathbf{u} =\mathbf{0}$, so $\mathcal{U}$ is invariant under $T_2,\ldots,T_m$.\\
Their restrictions to $\mathcal{U}$ still commute, and by induction have a common eigenvector $\mathbf{v}\in\mathcal{U}$, $\mathbf{v}\neq\mathbf{0}$.\\
Since $\mathbf{v}\in E_{\lambda}(T_1)$, it is also an eigenvector of $T_1$. For general $\mathcal{F}$, the subspace $\operatorname{span}(\mathcal{F})$ of $\LL(\mathcal{V})$ is finite-dimensional. Choose a basis $T_1,\ldots,T_m$ of this span from the spanning set $\mathcal{F}$.\\
The finite case gives a nonzero $\mathbf{v}$ and scalars $\mu_j$ with $T_j\mathbf{v}=\mu_j\mathbf{v}$. If $S\in\mathcal{F}$, write $S=\sum_j a_jT_j$.\\
Then $S\mathbf{v}=(\sum_j a_j\mu_j)\mathbf{v}$, so $\mathbf{v}$ is an eigenvector of every member of $\mathcal{F}$.\\
\hlt{Structural Commentary}: Commutation makes one eigenspace invariant, reducing the family by one.\\
\hlt{Key Insight}: Find the common eigenvector inside an eigenspace.
\end{proof}
\end{theorem}

\begin{theorem} \label{thm:upper-triangular-criteria} \hlt{Criteria for an Upper-Triangular Operator Matrix}\\
Let $T\in\LL(\mathcal{V})$ and let $\beta=(\mathbf{v}_1,\ldots,\mathbf{v}_n)$ be a basis of $\mathcal{V}$. The following are equivalent:
\begin{enumerate}[label=(\roman*)]
\setlength{\itemsep}{0pt}
\item $[T]_{\beta\leftarrow\beta}$ is upper triangular;
\item $T\mathbf{v}_j\in\operatorname{span}(\mathbf{v}_1,\ldots,\mathbf{v}_j)$ for each $j=1,\ldots,n$;
\item $\operatorname{span}(\mathbf{v}_1,\ldots,\mathbf{v}_j)$ is invariant under $T$ for each $j=1,\ldots,n$.
\end{enumerate}
\begin{proof}
\begin{enumerate}[label=(\roman*).]
\setlength{\itemsep}{0pt}
\item[$(i)$]$\Rightarrow (ii)$
The $j$th column of $[T]_{\beta\leftarrow\beta}$ is the coordinate column of $T\mathbf{v}_j$ in the basis $\beta$. Since the matrix is upper triangular, this column has zero entries below position $j$, so $T\mathbf{v}_j\in\operatorname{span}(\mathbf{v}_1,\ldots,\mathbf{v}_j)$.
\item[$(ii)$]$\Rightarrow (i)$
If $T\mathbf{v}_j\in\operatorname{span}(\mathbf{v}_1,\ldots,\mathbf{v}_j)$ for every $j$, then the $j$th coordinate column of $T\mathbf{v}_j$ has zero entries below position $j$. Hence $[T]_{\beta\leftarrow\beta}$ is upper triangular.
\item[$(ii)$]$\Rightarrow (iii)$
If $\mathbf{u}\in\operatorname{span}(\mathbf{v}_1,\ldots,\mathbf{v}_j)$, then linearity gives $T\mathbf{u}$ as a linear combination of $T\mathbf{v}_1,\ldots,T\mathbf{v}_j$, all of which lie in $\operatorname{span}(\mathbf{v}_1,\ldots,\mathbf{v}_j)$. Hence this span is invariant.
\item[$(iii)$]$\Rightarrow (ii)$
Apply invariance of $\operatorname{span}(\mathbf{v}_1,\ldots,\mathbf{v}_j)$ to $\mathbf{v}_j$.
\end{enumerate}
\hlt{Structural Commentary}: Upper triangularity records an invariant flag.\\
\hlt{Key Insight}: Columns below the diagonal vanish exactly when $T\mathbf{v}_j$ uses only earlier basis vectors and itself.
\end{proof}
\end{theorem}

\begin{corollary} \label{cor:upper-triangular-orthonormal-basis} \hlt{Upper-Triangular Form in an Orthonormal Basis}\\
Let $\mathcal{V}$ be a finite-dimensional inner-product space and let $T\in\LL(\mathcal{V})$. If $T$ has an upper-triangular matrix with respect to some basis of $\mathcal{V}$, then $T$ has an upper-triangular matrix with respect to some orthonormal basis of $\mathcal{V}$.
\begin{proof}
Let $\beta=(\mathbf{v}_1,\ldots,\mathbf{v}_n)$ be a basis for which $[T]_{\beta\leftarrow\beta}$ is upper triangular. By Theorem \ref{thm:upper-triangular-criteria}, each subspace
\begin{align}
\mathcal{U}_j=\operatorname{span}(\mathbf{v}_1,\ldots,\mathbf{v}_j),\qquad j=1,\ldots,n,
\nonumber
\end{align}
is invariant under $T$. Apply Gram--Schmidt to $\beta$ and obtain an orthonormal basis $\gamma=(\mathbf{e}_1,\ldots,\mathbf{e}_n)$ with
\begin{align}
\operatorname{span}(\mathbf{e}_1,\ldots,\mathbf{e}_j)=\operatorname{span}(\mathbf{v}_1,\ldots,\mathbf{v}_j)=\mathcal{U}_j.
\nonumber
\end{align}
Thus each $\operatorname{span}(\mathbf{e}_1,\ldots,\mathbf{e}_j)$ is invariant under $T$, and Theorem \ref{thm:upper-triangular-criteria} shows that $[T]_{\gamma\leftarrow\gamma}$ is upper triangular.\\
\hlt{Structural Commentary}: Upper triangularity is preservation of a nested flag of subspaces. Gram--Schmidt preserves that flag while making the basis orthonormal.\\
\hlt{Key Insight}: Triangular form depends on the flag, not on the particular basis used inside the flag.
\end{proof}
\end{corollary}

\begin{theorem} \label{thm:complex-upper-triangularization} \hlt{Complex Upper Triangularisation}\\
Let $\mathcal{V}$ be a finite-dimensional complex vector space and let $T\in\LL(\mathcal{V})$. Then $T$ has an upper-triangular matrix with respect to some basis of $\mathcal{V}$.
\begin{proof} Induct on $n=\dim(\mathcal{V})$. If $n\leq1$, the result is immediate.\\
Suppose $n>1$ and the result holds in smaller dimensions. Choose an eigenvector $\mathbf{v}_1$ of $T$ by Theorem \ref{thm:complex-operator-has-eigenvalue}, and set $\mathcal{U}=\operatorname{span}(\mathbf{v}_1)$. Then $\mathcal{U}$ is $T$-invariant, so the quotient operator $T/\mathcal{U}$ is defined on $\mathcal{V}/\mathcal{U}$, which has dimension $n-1$. By induction, choose a basis $(\mathbf{v}_2+\mathcal{U},\ldots,\mathbf{v}_n+\mathcal{U})$ of $\mathcal{V}/\mathcal{U}$ with respect to which $T/\mathcal{U}$ is upper triangular, and choose representatives $\mathbf{v}_2,\ldots,\mathbf{v}_n\in\mathcal{V}$.\\
The lifted list $(\mathbf{v}_1,\ldots,\mathbf{v}_n)$ is a basis of $\mathcal{V}$: any relation descends to the quotient basis, forcing all coefficients of $\mathbf{v}_2,\ldots,\mathbf{v}_n$ to be zero, and then the coefficient of $\mathbf{v}_1$ is zero. Now $T\mathbf{v}_1\in\operatorname{span}(\mathbf{v}_1)$.\\
For $j\geq2$, upper triangularity of $T/\mathcal{U}$ gives $T\mathbf{v}_j+\mathcal{U} =(T/\mathcal{U})(\mathbf{v}_j+\mathcal{U}) \in\operatorname{span}(\mathbf{v}_2+\mathcal{U},\ldots,\mathbf{v}_j+\mathcal{U})$, so $T\mathbf{v}_j\in\operatorname{span}(\mathbf{v}_1,\ldots,\mathbf{v}_j)$. Hence Theorem \ref{thm:upper-triangular-criteria} gives an upper-triangular matrix for $T$ in the basis $(\mathbf{v}_1,\ldots,\mathbf{v}_n)$.\\
\hlt{Structural Commentary}: Lift an upper-triangular quotient basis above an eigenline.\\
\hlt{Key Insight}: The quotient builds the remaining invariant flag.
\end{proof}
\end{theorem}

\begin{definition} \label{def:almost-upper-triangular-real-matrix} \hlt{Almost Upper-Triangular Real Matrices}\\
Let $A\in\operatorname{Mat}(n,n,\mathbb{R})$. The matrix $A$ is \hlt{almost upper triangular} if there exist $d_1,\ldots,d_m\in\{1,2\}$ with $\sum_jd_j=n$ such that, in the corresponding block form $A=[A_{ij}]_{i,j=1}^{m}$,
\begin{enumerate}[label=(\roman*)]
\setlength{\itemsep}{0pt}
\item $A_{ij}=0$ for $i>j$;
\item each diagonal block is either $A_{jj}=[\lambda_j]$, $\lambda_j\in\mathbb{R}$, or $A_{jj}=
\begin{bmatrix}
a_j&-b_j\\
b_j&a_j
\end{bmatrix}$, $a_j,b_j\in\mathbb{R},\ b_j\neq0$.
\end{enumerate}
The matrix $A$ is in \hlt{real Schur block form} if only (i) holds for some such block sizes $d_j\in\{1,2\}$.
\end{definition}

\begin{theorem} \label{thm:real-almost-upper-triangularization} \hlt{Real Almost Upper Triangularisation}\\
Let $\dim_{\mathbb{R}}\mathcal{V}<\infty$ and let $T\in\LL(\mathcal{V})$. Then $\exists$ a basis $\beta$ of $\mathcal{V}$ such that $[T]_{\beta\leftarrow\beta}$ is almost upper triangular.
\begin{proof}
Induct on $n=\dim_{\mathbb{R}}\mathcal{V}$. If $n\leq1$, clear.\\
For $n\geq2$, Theorem \ref{thm:complex-operator-has-eigenvalue} gives an eigenpair $T_{\mathbb{C}}(\mathbf{u}+i\mathbf{v})=(a+bi)(\mathbf{u}+i\mathbf{v})$, where $\mathbf{u}+i\mathbf{v}\neq\mathbf{0}$.\\
Thus $T\mathbf{u}=a\mathbf{u}-b\mathbf{v}$ and $T\mathbf{v}=b\mathbf{u}+a\mathbf{v}$.\\
If $b=0$, choose $\mathbf{0}\neq\mathbf{x}\in\{\mathbf{u},\mathbf{v}\}$, set $\mathcal{U}=\operatorname{span}(\mathbf{x})$, $\alpha=(\mathbf{x})$, so $[T|_{\mathcal{U}}]_{\alpha\leftarrow\alpha}=[a]$.\\
If $b\neq0$, set $\mathcal{U}=\operatorname{span}(\mathbf{u},\mathbf{v})$, $\alpha=(\mathbf{u},-\mathbf{v})$; then $\mathbf{u},\mathbf{v}$ are independent and $[T|_{\mathcal{U}}]_{\alpha\leftarrow\alpha}=
\begin{bmatrix}
a&-b\\
b&a
\end{bmatrix}$.\\
In either case, $T(\mathcal{U})\subseteq\mathcal{U}$, $\dim\mathcal{U}\in\{1,2\}$, and the first diagonal block is admissible.\\
By induction, $T/\mathcal{U}$ has an almost upper-triangular matrix in some quotient basis $\overline{\gamma}$.\\
Lift $\overline{\gamma}$ to $\gamma\subseteq\mathcal{V}$ and set $\beta=(\alpha,\gamma)$. Then $[T]_{\beta\leftarrow\beta}$ is block upper triangular with first block $[T|_{\mathcal{U}}]_{\alpha\leftarrow\alpha}$ and quotient block $[T/\mathcal{U}]_{\overline{\gamma}\leftarrow\overline{\gamma}}$, hence is almost upper triangular.\\
\hlt{Structural Commentary}: Real triangularisation starts with an invariant line or plane.\\
\hlt{Key Insight}: Complex eigenvectors split into real and imaginary parts; those parts form the first real block.
\end{proof}
\end{theorem}

\begin{corollary} \label{cor:real-schur-block-form-orthonormal} \hlt{Orthonormal Real Schur Block Form}\\
Let $\dim_{\mathbb{R}}\mathcal{V}<\infty$ and let $T\in\LL(\mathcal{V})$. Then $\exists$ an orthonormal basis $\gamma$ of $\mathcal{V}$ s.t. $[T]_{\gamma\leftarrow\gamma}$ is in real Schur block form.
\begin{proof}
By Theorem \ref{thm:real-almost-upper-triangularization}, choose $\beta=(\mathbf{v}_1,\ldots,\mathbf{v}_n)$ with block sizes $d_j\in\{1,2\}$.\\
Put $s_k=d_1+\cdots+d_k$ and $\mathcal{U}_k=\operatorname{span}(\mathbf{v}_1,\ldots,\mathbf{v}_{s_k})$. Then $T(\mathcal{U}_k)\subseteq\mathcal{U}_k$ for all $k$.\\
Apply Gram--Schmidt to $\beta$ and obtain an orthonormal basis $\gamma=(\mathbf{e}_1,\ldots,\mathbf{e}_n)$ with $\operatorname{span}(\mathbf{e}_1,\ldots,\mathbf{e}_{s_k})=\mathcal{U}_k$, $1\leq k\leq m$. Thus $[T]_{\gamma\leftarrow\gamma}$ is block upper triangular for the same $d_j\in\{1,2\}$, hence is in real Schur block form.\\
\hlt{Structural Commentary}: Gram--Schmidt preserves the invariant block flag.\\
\hlt{Key Insight}: Orthonormalization changes blocks, not the flag.
\end{proof}
\end{corollary}

\begin{proposition} \label{prop:upper-triangular-invertibility-criterion} \hlt{Invertibility from an Upper-Triangular Matrix}\\
Suppose $T\in\LL(\mathcal{V})$ has an upper-triangular matrix with diagonal entries $\lambda_1,\ldots,\lambda_n$ with respect to some basis of $\mathcal{V}$. Then $T$ is invertible if and only if $\lambda_j\neq0$ for every $j$.
\begin{proof}
Let $\beta=(\mathbf{v}_1,\ldots,\mathbf{v}_n)$ be such a basis.
\begin{enumerate}[label=\roman*.]
\setlength{\itemsep}{0pt}
\item[$(\Leftarrow)$]
If every $\lambda_j\neq0$, then upper-triangular matrix $[T]_{\beta\leftarrow\beta}$ is invertible by Theorem \ref{thm:inverse-upper-triangular}, so $T$ is invertible.
\item[$(\Rightarrow)$]
Suppose $T$ is invertible. If $\lambda_j=0$ for some $j$, set $\mathcal{U}_j=\operatorname{span}(\mathbf{v}_1,\ldots,\mathbf{v}_j)$ and $\mathcal{U}_0=\{\mathbf{0}\}$. Upper triangularity gives $T(\mathcal{U}_j)\subseteq\mathcal{U}_{j-1}$ because the coefficient of $\mathbf{v}_j$ in $T\mathbf{v}_j$ is $\lambda_j=0$ and $T\mathbf{v}_i\in\mathcal{U}_{j-1}$ for $i<j$.\\
Since $\dim(\mathcal{U}_j)>\dim(\mathcal{U}_{j-1})$, the map $\mathcal{U}_j\to\mathcal{U}_{j-1}$ induced by $T$ is not injective. Thus there is a nonzero $\mathbf{u}\in\mathcal{U}_j$ with $T\mathbf{u}=\mathbf{0}$, contradicting invertibility. Hence every $\lambda_j\neq0$.
\end{enumerate}
\hlt{Structural Commentary}: Diagonal entries control each flag layer.\\
\hlt{Key Insight}: A zero diagonal entry forces a kernel vector.
\end{proof}
\end{proposition}

\begin{proposition} \label{prop:eigenvalues-of-upper-triangular-operator} \hlt{Eigenvalues from an Upper-Triangular Matrix}\\
Suppose $T\in\LL(\mathcal{V})$ has an upper-triangular matrix with diagonal entries $\lambda_1,\ldots,\lambda_n$ with respect to some basis of $\mathcal{V}$. Then the eigenvalues of $T$ are precisely $\lambda_1,\ldots,\lambda_n$.
\begin{proof}
For $\lambda\in\mathbb{F}$, the matrix of $T-\lambda I_{\mathcal{V}}$ in the same basis is upper triangular with diagonal entries $\lambda_1-\lambda,\ldots,\lambda_n-\lambda$. By Proposition \ref{prop:upper-triangular-invertibility-criterion}, $T-\lambda I_{\mathcal{V}}$ is not invertible if and only if $\lambda=\lambda_j$ for some $j$. By Theorem \ref{thm:equivalent-eigenvalue-conditions}, this is equivalent to $\lambda$ being an eigenvalue of $T$.\\
\hlt{Structural Commentary}: Scalar shifts preserve triangular form.\\
\hlt{Key Insight}: Eigenvalues occur exactly when one shifted diagonal entry becomes zero.
\end{proof}
\end{proposition}

\begin{perspective} \label{per:upper-triangular-viewpoint} \hlt{Why Triangular Form Matters}\\
Upper-triangular form is the first useful simplification beyond finding one eigenvector. It preserves enough of the operator to read eigenvalues from the diagonal while still allowing nontrivial coupling above the diagonal.
\end{perspective}

\begin{definition} \label{def:diagonalizable-operator} \hlt{Diagonalizable Operator}\\
An operator $T\in\LL(\mathcal{V})$ is \hlt{diagonalizable} if $T$ has a diagonal matrix with respect to some basis of $\mathcal{V}$.
\end{definition}

\begin{proposition} \label{prop:eigenspaces-for-distinct-eigenvalues-form-direct-sum} \hlt{Eigenspaces for Distinct Eigenvalues Form a Direct Sum}\\
Let $\mathcal{V}$ be finite-dimensional, let $T\in\LL(\mathcal{V})$, and suppose $\lambda_1,\ldots,\lambda_m$ are distinct eigenvalues of $T$.\\
Then $E_{\lambda_1}(T)+\cdots+E_{\lambda_m}(T)$ is a direct sum. Consequently, $\dim E_{\lambda_1}(T)+\cdots+\dim E_{\lambda_m}(T)\leq\dim(\mathcal{V})$.
\begin{proof}
Suppose $\mathbf{u}_1+\cdots+\mathbf{u}_m=\mathbf{0}$ with $\mathbf{u}_j\in E_{\lambda_j}(T)$. If some $\mathbf{u}_j$ were nonzero, then the nonzero vectors among $\mathbf{u}_1,\ldots,\mathbf{u}_m$ would be eigenvectors corresponding to distinct eigenvalues, and Theorem \ref{thm:distinct-eigenvalues-independent-eigenvectors} would make them linearly independent. This contradicts their sum being $\mathbf{0}$. Hence every $\mathbf{u}_j=\mathbf{0}$, so the sum is direct.\\
By Theorem \ref{thm:direct-sum-iff-dimensions-add}, $\dim(E_{\lambda_1}(T)+\cdots+E_{\lambda_m}(T))=\dim E_{\lambda_1}(T)+\cdots+\dim E_{\lambda_m}(T)$.\\
The left side is the dimension of a subspace of $\mathcal{V}$, so Proposition \ref{prop:dimension-bounds-and-basis-criteria} gives the stated inequality.\\
\hlt{Structural Commentary}: Distinct eigenspaces form independent pieces.\\
\hlt{Key Insight}: A zero sum across them would contradict eigenvector independence.
\end{proof}
\end{proposition}

\begin{theorem} \label{thm:conditions-equivalent-to-diagonalizability} \hlt{Conditions Equivalent to Diagonalizability}\\
Let $\mathcal{V}$ be finite-dimensional, let $T\in\LL(\mathcal{V})$, and let $\lambda_1,\ldots,\lambda_m$ denote the distinct eigenvalues of $T$. Set $n=\dim(\mathcal{V})$. The following are equivalent:
\begin{enumerate}[label=(\roman*)]
\setlength{\itemsep}{0pt}
\item $T$ is diagonalizable;
\item $\mathcal{V}$ has a basis consisting of eigenvectors of $T$;
\item $\exists$ one-dimensional subspaces $\mathcal{U}_1,\ldots,\mathcal{U}_n$ of $\mathcal{V}$, each invariant under $T$, such that $\mathcal{V}=\mathcal{U}_1\oplus\cdots\oplus\mathcal{U}_n$;
\item $\mathcal{V}=E_{\lambda_1}(T)\oplus\cdots\oplus E_{\lambda_m}(T)$;
\item $\dim(\mathcal{V})=\dim E_{\lambda_1}(T)+\cdots+\dim E_{\lambda_m}(T)$.
\end{enumerate}
\begin{proof}
\begin{enumerate}[label=(\roman*).]
\setlength{\itemsep}{0pt}
\item[$(i)$]$\Rightarrow (ii)$ If $T$ is diagonalizable, then $T$ has a diagonal matrix with respect to some basis $(\mathbf{v}_1,\ldots,\mathbf{v}_n)$. Hence $T\mathbf{v}_j=a_j\mathbf{v}_j$ for each $j$, so this basis consists of eigenvectors of $T$.
\item[$(ii)$]$\Rightarrow (i)$ If $(\mathbf{v}_1,\ldots,\mathbf{v}_n)$ is a basis of eigenvectors of $T$, then $T\mathbf{v}_j=a_j\mathbf{v}_j$ for each $j$. Thus the matrix of $T$ with respect to this basis is diagonal.
\item[$(ii)$]$\Rightarrow (iii)$ Let $(\mathbf{v}_1,\ldots,\mathbf{v}_n)$ be an eigenbasis, and put $\mathcal{U}_j=\operatorname{span}(\mathbf{v}_j)$. Each $\mathcal{U}_j$ is one-dimensional and invariant under $T$ by Proposition \ref{prop:one-dimensional-invariant-subspaces-and-eigenvectors}; the basis decomposition gives $\mathcal{V}=\mathcal{U}_1\oplus\cdots\oplus\mathcal{U}_n$.
\item[$(iii)$]$\Rightarrow (ii)$ Choose nonzero $\mathbf{v}_j\in\mathcal{U}_j$. Proposition \ref{prop:one-dimensional-invariant-subspaces-and-eigenvectors} makes each $\mathbf{v}_j$ an eigenvector of $T$, and the direct-sum decomposition makes $(\mathbf{v}_1,\ldots,\mathbf{v}_n)$ a basis of $\mathcal{V}$.
\item[$(ii)$]$\Rightarrow (iv)$ Every vector in an eigenbasis lies in one of $E_{\lambda_1}(T),\ldots,E_{\lambda_m}(T)$, so these eigenspaces span $\mathcal{V}$. Proposition \ref{prop:eigenspaces-for-distinct-eigenvalues-form-direct-sum} shows their sum is direct.
\item[$(iv)$]$\Rightarrow (v)$ If $\mathcal{V}=E_{\lambda_1}(T)\oplus\cdots\oplus E_{\lambda_m}(T)$, then dimensions add by Theorem \ref{thm:direct-sum-iff-dimensions-add}.
\item[$(v)$]$\Rightarrow (ii)$ For each $j$, choose a basis $\beta_j$ of $E_{\lambda_j}(T)$, and concatenate these bases into a list $\beta$. The length of $\beta$ is $\dim(\mathcal{V})$. If a linear combination of vectors in $\beta$ is zero, grouping terms by eigenspace gives $\mathbf{u}_1+\cdots+\mathbf{u}_m=\mathbf{0}$ with $\mathbf{u}_j\in E_{\lambda_j}(T)$. Proposition \ref{prop:eigenspaces-for-distinct-eigenvalues-form-direct-sum} gives $\mathbf{u}_j=\mathbf{0}$ for every $j$; since each $\beta_j$ is independent, all coefficients are zero. Thus $\beta$ is independent, Proposition \ref{prop:full-dimension-basis-criteria} makes $\beta$ a basis of $\mathcal{V}$ consisting of eigenvectors of $T$.
\end{enumerate}
\hlt{Structural Commentary}: Diagonalizability means the eigenspaces exhaust $\mathcal{V}$. Directness is automatic for distinct eigenspaces, so the real test is spanning, or equivalently the dimension count.\\
\hlt{Key Insight}: A diagonal basis is exactly an eigenbasis.
\end{proof}
\end{theorem}
